\documentclass[12pt]{article}
\usepackage{amsmath,amssymb}

\begin{document}
\section*{{2018 AIME II Problems/Problem 14}}
Source: AoPS Wiki

\subsection*{Solution}
Let the sides $\overline{AB}$ and $\overline{AC}$ be tangent to $\omega$ at $Z$ and $W$ , respectively. Let $\alpha = \angle BAX$ and $\beta = \angle AXC$ . Because $\overline{PQ}$ and $\overline{BC}$ are both tangent to $\omega$ and $\angle YXC$ and $\angle QYX$ subtend the same arc of $\omega$ , it follows that $\angle AYP = \angle QYX = \angle YXC = \beta$ . By equal tangents, $PZ = PY$ . Applying the Law of Sines to $\triangle APY$ yields \[\frac{AZ}{AP} = 1 + \frac{ZP}{AP} = 1 + \frac{PY}{AP} = 1 + \frac{\sin\alpha}{\sin\beta}.\] Similarly, applying the Law of Sines to $\triangle ABX$ gives \[\frac{AZ}{AB} = 1 - \frac{BZ}{AB} = 1 - \frac{BX}{AB} = 1 - \frac{\sin\alpha}{\sin\beta}.\] It follows that \[2 = \frac{AZ}{AP} + \frac{AZ}{AB} = \frac{AZ}3 + \frac{AZ}7,\] implying $AZ = \tfrac{21}5$ . Applying the same argument to $\triangle AQY$ yields \[2 = \frac{AW}{AQ} + \frac{AW}{AC} = \frac{AZ}{AQ} + \frac{AZ}{AC} = \frac{21}5\left(\frac{1}{AQ} + \frac 18\right),\] from which $AQ = \tfrac{168}{59}$ . The requested sum is $168 + 59 = \boxed{227}$ .

\end{document}
